% English manuscript framework, grounded in the current repository.
% Compile with: pdflatex knowledge_discover_benchmark_outline.tex (twice).
% Self-contained: external figures and a bibliography database are not required.
% Evidence: run_benchmark.py, docs/benchmark_metrics.md,
% configs/benchmark/, kd/dataset/_catalog.py, kd/dataset/_registry.py.
% This is a writing framework, not a report of completed benchmark experiments.
\documentclass[11pt,a4paper]{article}
\usepackage[margin=1in]{geometry}
\usepackage[T1]{fontenc}
\usepackage{lmodern}
\usepackage{amsmath,amssymb}
\usepackage{tabularx,array}
\IfFileExists{booktabs.sty}{\usepackage{booktabs}}{%
  \newcommand{\toprule}{\hline}
  \newcommand{\midrule}{\hline}
  \newcommand{\bottomrule}{\hline}}
\usepackage{xcolor}
\newcolumntype{Y}{>{\raggedright\arraybackslash}X}
\setlength{\emergencystretch}{2em}
\definecolor{draftblue}{RGB}{25,75,135}
\newcommand{\todo}[1]{\par\smallskip\noindent{\color{draftblue}\textbf{[TO COMPLETE]} #1}\par\smallskip}
\newcommand{\uncertain}[1]{\par\smallskip\noindent{\color{draftblue}\textbf{[TO VERIFY]} #1}\par\smallskip}
\newcommand{\planned}[1]{\par\smallskip\noindent{\color{draftblue}\textbf{[PROPOSED EXPERIMENT]} #1}\par\smallskip}
\newcommand{\pending}{\textit{TBD}}
\newcommand{\code}[1]{\texttt{\detokenize{#1}}}
\newcommand{\figureplaceholder}[2]{%
  \fbox{\begin{minipage}[c][#1][c]{0.94\linewidth}
    \centering\small\textbf{Figure to be supplied}\par\medskip
    \begin{minipage}{0.9\linewidth}\raggedright #2\end{minipage}
  \end{minipage}}}
\title{KnowledgeDiscover: A Unified Benchmark for Symbolic Regression\\
and Discovery of Ordinary and Partial Differential Equations\\
\large Working manuscript framework}
\author{\pending{}: Authors and affiliations}
\date{\pending{}: Submission date and venue}

\begin{document}
\maketitle
\begin{abstract}
Recovering interpretable mathematical relationships from observations requires
evaluation beyond predictive accuracy. A discovered expression may predict well
without recovering the underlying structure, while an apparently correct
differential equation may produce unstable forecasts. We present
KnowledgeDiscover, a unified evaluation framework spanning symbolic regression
(SR), ordinary differential equation (ODE) discovery, and partial differential
equation (PDE) discovery. The framework combines task-specific data adapters,
explicit holdout protocols, isolated execution, and measurements of prediction,
symbolic recovery, dynamical consistency, and computational cost. Its current
runner registers 13 baseline adapters with task-dependent applicability.
SR and ODE evaluation includes normalized prediction error, coefficient of
determination, expression complexity, and algebraic recovery when references
are available. PDE evaluation separates exact support recovery, coefficient
accuracy, held-out equation residuals, and multi-step rollout error.
\textbf{[RESULTS TO INSERT: final dataset and instance counts, evaluated methods,
number of seeds, and two or three quantitative findings with uncertainty.]}
\textbf{[CONCLUSION TO INSERT: an evidence-supported implication for selecting
equation-discovery methods.]}
\end{abstract}

\noindent\textbf{Keywords:} symbolic regression; equation discovery; scientific
machine learning; reproducible benchmarking; symbolic recovery.

\section{Introduction}
\label{sec:introduction}
\subsection{Motivation}
Describe the scientific objective: infer compact, interpretable relationships
from measurements of static responses, trajectories, or spatiotemporal fields.
Explain why test-set accuracy alone does not establish recovery of a governing
law, and why runtime, failures, and data requirements affect practical utility.
Use one concrete motivating example linking a static constitutive relation,
an ODE, and a PDE without implying that all methods solve all three tasks.

\todo{Insert the target application and audience. Add verified citations for
scientific equation discovery and the limitations of prediction-only evaluation.}

\subsection{Evaluation gaps and research questions}
Position the paper around the following questions; these are questions to be
tested, not conclusions already supported by experiments.
\begin{enumerate}
\renewcommand{\labelenumi}{RQ\arabic{enumi}.}
  \item How do predictive accuracy and symbolic recovery agree or disagree
  across SR, ODE, and PDE tasks?
  \item For PDEs, how are exact support recovery, coefficient error, held-out
  residuals, and rollout accuracy related?
  \item How do measurement noise, sampling density, and distribution shifts
  change the relative performance of applicable methods?
  \item What tradeoffs arise among accuracy, expression complexity, execution
  reliability, runtime, and memory use?
\end{enumerate}

\subsection{Contributions and scope}
Frame the intended contributions as (i) a shared execution and reporting
infrastructure with task-specific adapters; (ii) a transparent evaluation
protocol that separates numerical fit, structural recovery, and dynamical
behavior; and (iii) a reproducible comparative study with explicit coverage
and failure accounting.
\uncertain{Establish novelty through related-work comparison before using
claims such as ``first,'' ``largest,'' or ``comprehensive.'' The contribution
is presently a benchmark and integration effort, not a new discovery algorithm.
Do not claim the empirical contribution until the full experiments are complete.}

\section{Related Work}
\label{sec:related}
\subsection{Symbolic regression and scientific discovery}
Discuss genetic programming, reinforcement-learning-guided search,
transformer-based expression generation, and physics-constrained symbolic
regression. Distinguish architectural ideas from the specific implementations
evaluated here.
\subsection{ODE and PDE identification}
Discuss sparse library regression, derivative estimation, neural surrogate
methods, symbolic search, and learned evolution operators. Explain the distinct
roles of discovering equations and predicting trajectories or fields.
\subsection{Existing benchmarks}
Compare task coverage, data provenance, ground-truth availability, split rules,
metric definitions, compute accounting, and reproducibility artifacts.

\begin{table}[htbp]
\centering\small
\caption{Related-benchmark comparison template. Every external entry requires
a verified citation and protocol inspection; blanks are not negative findings.}
\label{tab:related}
\begin{tabularx}{\linewidth}{lYYYY}
\toprule
Benchmark & Tasks/data & Recovery definition & Holdout protocol & Cost/failures \\
\midrule
\pending{} benchmark A & \pending & \pending & \pending & \pending \\
\pending{} benchmark B & \pending & \pending & \pending & \pending \\
\pending{} benchmark C & \pending & \pending & \pending & \pending \\
KnowledgeDiscover & SR/ODE/PDE & Task-specific & See Sec.~\ref{sec:protocol} & Logged \\
\bottomrule
\end{tabularx}
\end{table}
\todo{Add verified primary references for each evaluated method, dataset, and
comparison benchmark. No external bibliographic entries have been invented for
this draft; complete the bibliography before submission.}

\section{Benchmark Design}
\label{sec:design}
\subsection{Problem formulation}
For SR, estimate an expression $\hat f$ from observations
$y_i=f^\star(\mathbf{x}_i)+\epsilon_i$.
For ODE discovery, estimate each component of
\begin{equation}
  \dot{\mathbf{z}}(t)=\mathbf{f}^\star(\mathbf{z}(t),\mathbf{p}),
\end{equation}
using state and parameter features and derivative targets constructed by the
adapter. For PDE discovery, consider equations representable as
\begin{equation}
  \partial_t u=F^\star(u,\nabla u,\nabla^2u,\ldots,\mathbf{x}).
\end{equation}
These equations describe the evaluation setting, not universal support for all
ODEs, coupled PDE systems, geometries, or boundary conditions. Multiple scalar
targets may produce multiple result rows within one execution case.

\subsection{Data inventory and reference fidelity}
The runner draws from the unified dataset registry, a symbolic-expression CSV
catalog, and a wave-breaking adapter. Candidate sources include synthetic SR
families, material-response regression, ball-drop trajectories, packaged PDE
fields, TLC time-series data, wake data, viscous-gravity-current data, and a
generated PDE family. Registry membership does not establish successful loading,
method compatibility, or inclusion in the final study.
Classify references as coefficient-bearing equations, structure-only equations,
or unavailable ground truth. Keep experimental observations separate from
simulations and generated expressions.

\begin{table}[htbp]
\centering\small
\caption{Dataset inventory to finalize from the frozen experiment manifest.
Counts refer to distinct tasks/instances, not simply registered names.}
\label{tab:datasets}
\begin{tabularx}{\linewidth}{>{\raggedright\arraybackslash}p{0.27\linewidth}YYYY}
\toprule
Task/source examples & Instances; targets & Size; domain & Reference & Split; provenance \\
\midrule
SR: Koza, Keijzer, other CSV families & \pending & \pending & Expressions & Provided split; \pending \\
SR: rubber, CYT, solid-response data & \pending & \pending & Verify per dataset & Official, group, or random \\
ODE: ball drop & \pending & \pending & Verify component equations & Whole trajectories \\
PDE: Burgers, KdV, Chafee--Infante & \pending & \pending & Coefficient-bearing & Temporal block \\
PDE: TLC, wake, wave breaking, VGS & \pending & \pending & Verify per dataset & Temporal block \\
PDE: generated sinus family & \pending & \pending & Verify metadata & Temporal block \\
\bottomrule
\end{tabularx}
\end{table}
\todo{For every included dataset, specify source citation, license, checksums,
units, dimensions, spatial and temporal resolution, sample/trajectory counts,
noise characteristics, boundary/initial conditions, and preprocessing. Deduplicate
aliases and related instances. The rubber train/test pair is one experiment.}
\uncertain{The catalog explicitly excludes a broken advection--diffusion loader
and unsupported steady-state TLC files. Audit generated-instance seeds and
independence across splits; finalize all counts after loading and eligibility checks.}

\subsection{Baseline adapters and applicability}
Table~\ref{tab:methods} describes the current runner, not completed experiment
coverage. Each adapter's operators, features, initialization, units, search
budget, and stopping rule must be documented for the publication.

\begin{table}[htbp]
\centering\small
\caption{Registered baseline adapters and implementation scope. GPU support
refers to supported components, not necessarily the entire search.}
\label{tab:methods}
\begin{tabularx}{\linewidth}{llY}
\toprule
Adapter ID & Registered tasks & Implementation note \\
\midrule
\code{dscv} & SR, ODE, PDE & DISCOVER adapter; current device interface is CPU. \\
\code{spr} & PDE & DISCOVER/SPR integration; PINN component supports CUDA. \\
\code{sga} & PDE & SGA integration; CPU in current runner. \\
\code{dlga} & PDE & Neural derivatives support CUDA; genetic search includes CPU work. \\
\code{deepmod} & PDE & DeepMoD neural and sparse-model components. \\
\code{pdefind} & PDE & PDE-FIND adapter using numerical regression. \\
\code{pdenet} & PDE & PDE-Net adapter with a model-based evolution predictor. \\
\code{weakform} & PDE & Independent scalar-1D WSINDy-PDE adaptation using compact
test functions, convolutional weak derivatives, scaling, and MSTLS. \\
\code{eqgpt} & PDE & EqGPT adapter with CUDA support. \\
\code{dso} & SR, ODE & PyTorch-based DSO integration; document departure from upstream. \\
\code{symbolicgpt} & SR, ODE & Transformer-based symbolic-regression adapter. \\
\code{gplearn} & SR, ODE & CPU genetic-programming regressor. \\
\code{physo} & SR, ODE & PhySO integration with selectable device. \\
\bottomrule
\end{tabularx}
\end{table}
\uncertain{SISSO++ v1.2.6 was installed and smoke-tested separately, but is not
registered in the inspected runner. Its inclusion requires an adapter, a baseline
configuration, applicable task definitions, and evaluation tests. Keep it out of
the reported 13 adapters until that work is complete. PySINDy installation alone
likewise does not establish a separately evaluated baseline.}
\todo{Resolve remaining display names and upstream citations. State the scalar-1D
scope of the WSINDy-PDE adaptation and distinguish it from upstream multidimensional,
multifield, trigonometric, automatic-support, and RGLS options. Audit pretrained
checkpoints, training corpora, and test-set overlap for all
learned priors. Specify which physics or dimensional information each method receives.}

\subsection{Execution and reproducibility infrastructure}
The runner builds model--dataset--seed cases, imports methods on demand, executes
cases in separate worker processes and working directories, and records resolved
parameters and devices. It supports timeout handling, partial results, and
configuration-aware resumption. Current scheduling is sequential at the case
level. Status values distinguish successful execution, errors, timeouts, and skips.
Artifacts include a manifest, case configurations and logs, raw result rows,
and task--method aggregates.

\begin{figure}[htbp]
\centering
\figureplaceholder{4cm}{Draw the actual workflow: dataset catalog and references
$\rightarrow$ task adapters and splits $\rightarrow$ model--dataset--seed manifest
$\rightarrow$ isolated workers $\rightarrow$ expression/prediction normalization
$\rightarrow$ task metrics and resource measurements $\rightarrow$ raw and aggregate
reports. Show missing-reference and unsupported-adapter branches explicitly.}
\caption{Planned architecture figure for the benchmark.}
\label{fig:architecture}
\end{figure}

\section{Evaluation Protocol}
\label{sec:protocol}
\subsection{Train/test separation}
\paragraph{SR.} Synthetic expression benchmarks use their provided benchmark
splits. Rubber data use their official paired split. Other regression datasets
use group holdout when group IDs exist and random holdout otherwise. Training
subsampling is seed-controlled and applies only to training examples.
\paragraph{ODE.} Entire trajectories are assigned to train or test before
conversion to regression arrays. First-order derivative targets are generated
separately for each subset, and boundary derivative estimates are discarded.
Current ODE scores concern derivative prediction, not integrated trajectory error.
\paragraph{PDE.} For $n_t\geq5$ frames and holdout fraction $\rho$, the current
adapter sets
\begin{equation}
 n_{\mathrm{test}}=\min\{n_t-3,\max\{2,\operatorname{round}(\rho n_t)\}\},
 \qquad k=n_t-n_{\mathrm{test}}.
\end{equation}
With zero-based indexing, training uses frames $0,\ldots,k-1$ and testing uses
$k,\ldots,n_t-1$. The evaluation view includes frame $k-1$ as a rollout initial
condition, but this frame is excluded from test-error scoring.

\begin{figure}[htbp]
\centering
\figureplaceholder{3.4cm}{Three panels: (a) SR official/group holdout; (b) complete
ODE trajectories in disjoint subsets; (c) PDE training prefix, final training
frame used as the initial condition, and future test block. Label derivative
stencil boundaries and distinguish observed test fields for residual evaluation
from predicted states for rollout evaluation.}
\caption{Planned split-protocol figure.}
\label{fig:splits}
\end{figure}
\todo{Audit all adapters for preprocessing fitted on training data only,
derivative stencils, normalization, and internal validation. Establish a validation
split or training-only tuning protocol before selecting publication hyperparameters.
Do not select configurations using the held-out test scores.}

\subsection{SR and ODE metrics}
Let $\sigma_y^2=N^{-1}\sum_i(y_i-\bar y)^2$ on the test targets. Report
\begin{equation}
 \mathrm{NRMSE}=\frac{\sqrt{N^{-1}\sum_i(\hat y_i-y_i)^2}}{\sigma_y},
 \qquad R^2=1-\frac{\sum_i(\hat y_i-y_i)^2}{\sum_i(y_i-\bar y)^2}.
\end{equation}
For ODE tasks, $y$ is a derivative component. Constant targets require explicit
degenerate-case handling: the current NRMSE is zero for an exact fit and unavailable
for a nonzero error on a constant target; $R^2$ is unavailable for any constant target.
Expression complexity is the number of variable, numeric, function, and operator
tokens in the emitted RHS before algebraic simplification. It is a lexical proxy,
not a proven representation-invariant measure of minimal description length.

For parseable expressions with trustworthy ground truth, algebraic recovery is
\begin{equation}
 I_{\mathrm{alg}}=\mathbb{1}\{\operatorname{simplify}(\hat f-f^\star)=0\}.
\end{equation}
Unavailable references and parsing failures yield unavailable scores. The
recovery rate averages evaluable indicators and must be accompanied by their count.
\uncertain{Symbolic simplification depends on parser conventions, variable
mapping, domain assumptions, and floating-point constants. Document these and
any numerical-equivalence supplement; do not silently substitute approximate
numerical agreement for this algebraic criterion.}

\subsection{PDE structure and coefficient metrics}
Normalize supported equations to an explicit RHS. Expand additive terms,
combine equivalent term keys, and remove coefficients with magnitude at most
$10^{-12}$, matching the current implementation. Let $S^\star$ and $\hat S$
be reference and discovered supports. Report
\begin{align}
 I_{\mathrm{support}} &= \mathbb{1}\{\hat S=S^\star\}, \\
 J &= \frac{|\hat S\cap S^\star|}{|\hat S\cup S^\star|}, &
 P &= \frac{|\hat S\cap S^\star|}{|\hat S|}, &
 R &= \frac{|\hat S\cap S^\star|}{|S^\star|}.
\end{align}
The implementation assigns $J=1$ when both supports are empty, and zero precision
or recall when their respective denominators are empty.
The output field \code{exact_symbolic_recovery} corresponds to
$I_{\mathrm{support}}$. In the paper, call it \emph{exact symbolic support recovery}
to distinguish it from coefficient-exact algebraic recovery. It ignores coefficient
magnitudes and signs after support extraction; even an incorrect coefficient sign
can therefore coexist with perfect support recovery.

When reliable numerical coefficients are available, align coefficients on
$\hat S\cup S^\star$, insert zeros for absent terms, and compute
\begin{equation}
 E_c=\frac{\|\hat{\mathbf c}-\mathbf c^\star\|_2}{\|\mathbf c^\star\|_2}.
\end{equation}
Do not compute $E_c$ from structure-only prefix strings. Coefficient-bearing
references currently exist in the runner for Burgers, KdV, and Chafee--Infante.
\todo{Verify those references against the precise packaged data generators,
including signs, units, nondimensionalization, and coordinate transformations.
State the support threshold and test its sensitivity before publication.}

\subsection{PDE equation residuals and multi-step rollout}
On observed held-out fields, evaluate
\begin{equation}
 r=\widehat{\partial_tu}-\hat F(u,\widehat{\nabla u},\ldots),
 \qquad E_{\mathrm{res}}=\frac{\sqrt{\operatorname{mean}(r^2)}}
 {\operatorname{std}(\widehat{\partial_tu})}.
\end{equation}
The evaluator uses grid-based finite differences, including repeated spatial
differentiation. A small residual measures agreement with those derivative
estimates; it is not independent evidence of physical correctness.

Initialize rollout at $\hat u_{k-1}=u_{k-1}$ and predict the complete held-out
time block without resetting to observed intermediate states. A compatible
model predictor is used when available; otherwise, the current fallback is
explicit Euler applied to a parseable RHS. Report rollout NRMSE, $R^2$, MSE,
step count, horizon, and integrator. Scoring pools held-out field values according
to the current evaluator.
\uncertain{Integrator stability and boundary treatment can dominate rollout
error, especially for stiff or dispersive PDEs. The current predictor-versus-Euler
comparison is not solver-controlled. Validate on the reference RHS and report
integrator-stratified results; a common converged solver is a proposed extension.
Unsupported or non-finite evaluations must be reported with their reasons.}

\subsection{Reliability, compute, and aggregation}
Execution success is distinct from scientific recovery. For a task--method group
with $N_{\mathrm{rows}}$ emitted rows, the current report computes
\begin{equation}
 p_{\mathrm{ok}}=\frac{N_{\mathrm{ok}}}{N_{\mathrm{rows}}},
 \qquad p_{\mathrm{timeout}}=\frac{N_{\mathrm{timeout}}}{N_{\mathrm{rows}}}.
\end{equation}
Errors and skips remain in these denominators. Runtime records training and
evaluation wall time at the instance/target level; case timeouts act at the worker
level. Peak CPU memory measures sampled worker-process-tree RSS when possible;
GPU memory measures the CUDA allocator peak, not total device memory.

Current metric means use available numeric values. Consequently, coverage and
evaluable counts are essential: unavailable values are not zero errors, and an
\code{ok} row need not contain every scientific metric. Report raw-status counts,
metric-specific coverage, and conditional recovery rates together.
\planned{Add dataset-balanced summaries and dataset-level bootstrap intervals,
with seeds nested within datasets, after defining the eligible comparison cohort.
Supplement current row-level rates with case-level rates and an attempted-run
recovery rate that treats failed attempts on reference-known eligible tasks as
failures. Mark these as additional analyses, not existing runner outputs.
Treat timeouts as censored runtime observations rather than completed fits.}

\section{Experimental Setup}
\label{sec:setup}
\subsection{Study matrix and budgets}
Freeze dataset instances, seeds, applicability rules, operator libraries,
hyperparameters, and compute budgets before collecting final test results.
The default configuration currently selects the smoke profile, seed 0, CPU
devices, a 0.2 holdout fraction, and no global timeout. The full profile increases
method-specific budgets but does not claim reproduction of the original papers.

\begin{table}[htbp]
\centering\small
\caption{Publication experiment settings. Environment setup is known; final
benchmark settings require confirmation and must match the released manifest.}
\label{tab:setup}
\begin{tabularx}{\linewidth}{lY}
\toprule
Item & Value/status \\
\midrule
Dataset inventory and reference tier & \pending{}: Table~\ref{tab:datasets}, hashes, inclusion rules \\
Seeds and repeat counts & \pending{}: choose before final runs; current default is one seed \\
Holdout and validation & Default test fraction 0.2; publication validation protocol \pending \\
Per-method budgets / timeout & \pending{}: include sample/epoch/search limits and wall-clock cap \\
CPU resources & Prior container provision: 8 CPU quota; CPU model and final allocation \pending \\
Memory & Prior container provision: 60 GiB limit; final run allocation \pending \\
GPU & Prior validated environment: NVIDIA RTX 4090; per-method assignment \pending \\
Software & Ubuntu 22.04, Python 3.9.23, PyTorch 2.5.1+cu121,
CUDA Toolkit 12.1.1; final lockfile/commit \pending \\
Thread settings & OpenMP/OpenBLAS default to 1; audit other workers and thread pools \\
Checkpoint provenance & \pending{}: pretrained/from-scratch status and training-data overlap \\
\bottomrule
\end{tabularx}
\end{table}
\uncertain{Installation checks, CUDA kernel checks, and a tiny SISSO++ fit verify
the environment only. They are not accuracy or efficiency evidence for the final
benchmark. Reconfirm hardware and software on the actual experiment session.}

\subsection{Main study and controlled extensions}
\planned{Run the finalized clean-data study across all eligible task--method
pairs and multiple seeds. Then add a predeclared subset for noise, sample-density,
and generalization experiments. Select levels and repetitions after a pilot,
without consulting the final test scores.}
\todo{Specify a noise model, e.g., additive noise scaled by the training-field
standard deviation; state whether test observations remain clean. Define whether
noise precedes derivative estimation. Specify trajectory/initial-condition shifts
and spatial or temporal subsampling without destroying task identifiability.}

\section{Results and Analysis}
\label{sec:results}
\todo{This section intentionally contains no rankings or numerical claims.
Populate it only from a frozen full-run manifest and its logs. Replace every
TBD and report uncertainty and evaluable counts.}

\subsection{RQ1: Prediction and algebraic recovery in SR and ODE}
Discuss each task separately, distinguishing synthetic reference-known problems
from empirical regression and reference-unknown ODEs. Test whether error and
complexity improvements correspond to increased algebraic recovery.
\begin{table}[htbp]
\centering\small
\caption{SR/ODE result template. Expand to one row per applicable method and
task; attach an evaluable count to each metric and specify the aggregation unit.}
\label{tab:sr-ode}
\begin{tabularx}{\linewidth}{llYYYYY}
\toprule
Task & Method & NRMSE $\downarrow$ & $R^2$ $\uparrow$ & Complexity $\downarrow$ & Recovery $\uparrow$ & Evaluable / eligible \\
\midrule
SR & \pending & \pending & \pending & \pending & \pending & \pending \\
ODE & \pending & \pending & \pending & \pending & \pending & \pending \\
\bottomrule
\end{tabularx}
\end{table}

\subsection{RQ2: PDE structure, coefficients, and dynamics}
Analyze structural mistakes, coefficient errors given correct support, and
rollout failures separately. Do not equate exact support recovery with complete
equation recovery. Stratify rollout comparisons by integrator and ground-truth tier.
\begin{table}[htbp]
\centering\small
\caption{PDE result template. Use multiple panels or landscape layout for the
final table; report precision/recall and metric-specific counts in the supplement.}
\label{tab:pde-results}
\begin{tabularx}{\linewidth}{lYYYYY}
\toprule
Method & Residual NRMSE $\downarrow$ & Coefficient error $\downarrow$ & Support Jaccard $\uparrow$ & Exact support rate $\uparrow$ & Rollout NRMSE $\downarrow$ \\
\midrule
\pending{} method A & \pending & \pending & \pending & \pending & \pending \\
\pending{} method B & \pending & \pending & \pending & \pending & \pending \\
Reference RHS control & \pending & Not applicable & Not applicable & Not applicable & \pending \\
\bottomrule
\end{tabularx}
\end{table}

\begin{figure}[htbp]
\centering
\figureplaceholder{4cm}{Panels: exact support rate versus coefficient error;
held-out residual versus rollout error; representative true and discovered PDEs.
Label sample counts and integrators. Include examples of correct support with
wrong coefficients, low residual with unstable rollout, and full agreement only
if such cases occur in the measured results.}
\caption{Planned analysis of distinct PDE discovery outcomes.}
\label{fig:pde-outcomes}
\end{figure}

\subsection{RQ3: Robustness and generalization}
\begin{figure}[htbp]
\centering
\figureplaceholder{4cm}{Use noise level and training-data fraction as separate
horizontal axes. Plot NRMSE, exact recovery, and execution success with dataset-level
uncertainty. Add an initial-condition or out-of-domain panel only if that experiment
is implemented. Keep task and reference tiers separate.}
\caption{Planned robustness study; experiments and perturbation levels remain to be specified.}
\label{fig:robustness}
\end{figure}
\todo{Describe observed trends, exceptions, and derivative-estimation sensitivity.
Separate interpolation results from true extrapolation; temporal holdout alone
does not establish generalization to unseen governing equations.}

\subsection{RQ4: Reliability and computational cost}
\begin{table}[htbp]
\centering\small
\caption{Compute and reliability template, repeated by task and hardware cohort.
Include errors/skips and completed versus censored runtime counts in final reporting.}
\label{tab:cost}
\begin{tabularx}{\linewidth}{lYYYYY}
\toprule
Method & Success rate & Timeout rate & Runtime (s) & Peak RSS (MB) & Peak GPU (MB) \\
\midrule
\pending{} method A & \pending & \pending & \pending & \pending & \pending \\
\pending{} method B & \pending & \pending & \pending & \pending & \pending \\
\bottomrule
\end{tabularx}
\end{table}
\begin{figure}[htbp]
\centering
\figureplaceholder{3.8cm}{Accuracy--runtime and accuracy--complexity Pareto plots
within matched task/hardware cohorts, plus stacked status proportions. Mark
timeouts and unsupported combinations; avoid plotting unavailable metrics as
zero. Do not aggregate SR, ODE, and PDE into an unexplained universal score.}
\caption{Planned cost, complexity, and reliability comparison.}
\label{fig:cost}
\end{figure}

\subsection{Ablations and qualitative case studies}
\planned{Prioritize three checks: (1) derivative-estimation and support-threshold
sensitivity; (2) reference-RHS and common-solver rollout controls; (3) compute-budget
sensitivity. Optionally quantify random-point versus trajectory/time-block holdout
as a leakage diagnostic, not as an alternative main ranking.}
\begin{table}[htbp]
\centering\small
\caption{Equation case-study template. Select examples by a documented rule
covering both successes and failures, with numerical metrics and units.}
\label{tab:equations}
\begin{tabularx}{\linewidth}{lYYYY}
\toprule
Dataset/method & Reference & Discovered equation & Scores & Interpretation \\
\midrule
\pending{} SR/ODE case & \pending & \pending & \pending & \pending \\
\pending{} PDE recovery & \pending & \pending & \pending & \pending \\
\pending{} PDE failure & \pending & \pending & \pending & \pending \\
\bottomrule
\end{tabularx}
\end{table}

\section{Discussion and Limitations}
Discuss what the results support about method selection for a given task,
data quality, reference availability, and compute budget. Address incomplete
adapter coverage, implementation differences from upstream methods, limited
ODE diversity, correlated dataset instances, missing physical units, and possible
pretraining contamination. Discuss derivative noise and numerical integration
as confounders, and parser-dependent recovery and lexical complexity as metric
limitations. Acknowledge that rollout evaluation for coupled fields and general
boundary conditions may require additional adapters.
\todo{Replace prospective discussion with claims tied to specific tables and
figures. State which limitations remain after the publication experiments. Include
the PhySO CUDA/multiprocessing constraint if it affects the measured configuration.}

\section{Conclusion}
Summarize the released evaluation framework and the empirical answers to the
research questions, using only completed evidence. Explain why prediction,
structural recovery, coefficient accuracy, and dynamical behavior should be
reported together. End with the most useful supported implication for scientific
users and a concrete remaining evaluation gap.
\todo{Insert two or three quantitative takeaways after the final analysis.
Do not infer benchmark superiority from environment or adapter smoke tests.}

\section*{Data and Code Availability}
\todo{Insert the repository/archive URL, release tag and commit, dataset licenses
and access instructions, environment lockfiles, container recipe or image digest,
raw manifests/logs/results, and scripts that generate all publication figures and
tables. Distinguish redistributable data from externally downloaded resources.}
\section*{Acknowledgments}
\todo{Authors, funding, computational resources, and upstream software/data credits.}
\section*{References}
\todo{Replace this section with a verified bibliography. Add primary sources
for all methods in Table~\ref{tab:methods}, every included dataset, related
benchmarks, and the numerical procedures discussed in the protocol.}

\appendix
\section{Complete Dataset and Eligibility Manifest}
Supply the full dataset--method matrix with load status, task dimensions,
response channels, reference tier, eligibility decision, and exclusion reason.
Report counts of datasets, instances, targets, seeds, scheduled cases, and emitted
rows separately. Preserve official splits and group/trajectory identities.

\section{Resolved Baseline Configurations}
Archive every baseline JSON plus resolved overrides. Record operator grammars,
candidate libraries, sparsity thresholds, neural architectures, optimization
settings, stopping rules, devices, and upstream/custom modifications. Document
physical priors, checkpoints, and any tuning search with its validation budget.

\section{Metric Implementation and Edge Cases}
Provide parser conventions, variable aliases, support threshold, coefficient
alignment, empty-support rules, constant-target handling, and missing-value
semantics. Document derivative stencils, boundary behavior, allowed RHS operators,
rollout initialization, integrators, time steps, and divergence handling. Explain
why \code{exact_recovery} and \code{exact_symbolic_recovery} have different meanings.

\section{Statistical Reporting and Failure Audit}
Publish per-dataset/per-seed values and metric-specific denominators alongside
aggregates. Separate unsupported inputs, loading failures, fitting failures,
parsing failures, metric failures, and timeouts. Describe bootstrap or other
uncertainty procedures if implemented, including resampling unit and seed.

\section{Reproduction and Completion Checklist}
\begin{enumerate}
  \item Freeze the final scientific scope, verified references, dataset licenses,
  and method/instance eligibility matrix.
  \item Finalize validation, seeds, compute limits, device assignments, and
  measurement conventions; archive the resolved manifest and environment.
  \item Verify task adapters, reference equations, numerical evaluators, and
  symbolic metric edge cases before running the main study.
  \item Run the main study and declared robustness/ablation experiments;
  preserve all failures and raw outputs.
  \item Generate the eight tables and five figures reserved in this framework
  from analysis scripts, adding confidence intervals and coverage counts.
  \item Replace all TO COMPLETE, TO VERIFY, PROPOSED EXPERIMENT, and TBD markers;
  reconcile the abstract, results, discussion, and conclusion with the evidence.
\end{enumerate}
\end{document}
